\section*{Capítulo \ref{ch:b1:faults-folds-fabrics} --- Fallas, pliegues y fábricas}

\begin{solution}{exo:b1:faults-folds-fabrics:1}
(a) Una \omterm{def:b1:faults-folds-fabrics:fault-types}{falla normal}. (b) Una \omterm{def:b1:faults-folds-fabrics:fault-types}{falla inversa}, casi de pura componente de \omterm{def:b1:geological-maps:attitude}{buzamiento}. (c) Una \omterm{def:b1:faults-folds-fabrics:fault-types}{falla de desgarre} \omterm{def:b1:faults-folds-fabrics:fault-types}{sinistral}.
(d) Una \omterm{def:b1:faults-folds-fabrics:fault-types}{falla de desgarre} \omterm{def:b1:faults-folds-fabrics:fault-types}{dextral}. (e) Una \omterm{def:b1:faults-folds-fabrics:fault-types}{falla oblicua}: el \omterm{def:b1:faults-folds-fabrics:fault-elements}{bloque superior} bajó y se desplazó a lo largo del \omterm{def:b1:geological-maps:attitude}{rumbo}, hacia la
izquierda; su componente de desgarre, $\cos\ang{40} = 0.77$ del \omterm{def:b1:faults-folds-fabrics:fault-elements}{deslizamiento neto}, es mayor que su componente de \omterm{def:b1:geological-maps:attitude}{buzamiento},
$0.64$, así que se llama falla \omterm{def:b1:faults-folds-fabrics:fault-types}{sinistral} con una componente normal. (f) Con un \omterm{def:b1:geological-maps:attitude}{buzamiento} de \ang{25}, menor que \ang{45}, la
\omterm{def:b1:faults-folds-fabrics:fault-types}{falla inversa} es un \omterm{def:b1:faults-folds-fabrics:fault-types}{cabalgamiento}.
\end{solution}

\begin{solution}{exo:b1:faults-folds-fabrics:2}
Con $|\sin r| = 1$, el salto es $300\sin\ang{55} = \qty{245.7}{m}$ y el \omterm{def:b1:cross-sections:throw}{salto horizontal}
$300\cos\ang{55} = \qty{172.1}{m}$.
\end{solution}

\begin{solution}{exo:b1:faults-folds-fabrics:3}
Para un \omterm{def:b1:faults-folds-fabrics:fold-classes}{pliegue erguido} cuyos \omterm{def:b1:faults-folds-fabrics:fold-elements}{flancos} buzan $\delta_1$ y $\delta_2$ en sentidos opuestos, $\iota = \ang{180} - \delta_1 - \delta_2$.
(a) \ang{110}: abierto. (b) \ang{35}: cerrado. (c) \ang{7}: isoclinal.
\end{solution}

\begin{solution}{exo:b1:faults-folds-fabrics:4}
Las \omterm{def:b1:faults-folds-fabrics:fault-types}{fallas normales} buzan $\ang{45} + \ang{35}/2 = \ang{62.5}$, y los \omterm{def:b1:faults-folds-fabrics:fault-types}{cabalgamientos}
$\ang{45} - \ang{17.5} = \ang{27.5}$. Las \omterm{def:b1:faults-folds-fabrics:fault-types}{fallas de desgarre} conjugadas forman cada una
\ang{27.5} con $\sigma_1$, a uno y otro lado de él, así que se cortan con \ang{55}.
\end{solution}

\begin{solution}{exo:b1:faults-folds-fabrics:5}
(a) La componente de desgarre es $400\cos\ang{30} = \qty{346.4}{m}$ según el \omterm{def:b1:geological-maps:attitude}{rumbo}, y la componente de \omterm{def:b1:geological-maps:attitude}{buzamiento}
$400\sin\ang{-30} = \qty{-200}{m}$, es decir, \qty{200}{m} hacia abajo según el \omterm{def:b1:geological-maps:attitude}{buzamiento}; el salto es $200\sin\ang{50} = \qty{153.2}{m}$ y
el \omterm{def:b1:cross-sections:throw}{salto horizontal} $200\cos\ang{50} = \qty{128.6}{m}$. (b) El \omterm{def:b1:geological-maps:attitude}{rumbo} es $000$, así que $\mathbf s = (0, 1, 0)$;
$\mathbf u = (-\cos\ang{50}, 0, \sin\ang{50}) =
(-0.643, 0, 0.766)$. Entonces
$\overrightarrow{AA'} = 400(0.866\,\mathbf s - 0.5\,\mathbf u) =
(128.6, 346.4, -153.2)$ en metros. Su tendencia es $\arctan(128.6/346.4) =
\ang{20.4}$ y su \omterm{def:b1:field-lab-methods-1:trend-plunge}{inmersión} $\arcsin(153.2/400) = \ang{22.5}$, que es también
$\arcsin(\sin\ang{30}\sin\ang{50})$. (c) Una \omterm{def:b1:faults-folds-fabrics:fault-types}{falla oblicua}: el \omterm{def:b1:faults-folds-fabrics:fault-elements}{cabeceo} queda
entre $0$ y \ang{-90}, más cerca de $0$: una falla \omterm{def:b1:faults-folds-fabrics:fault-types}{sinistral} con una
componente normal.
\end{solution}

\begin{solution}{exo:b1:faults-folds-fabrics:6}
(a) Las dos fallas forman \ang{60} entre sí; $\sigma_1$ biseca el ángulo agudo, así que es horizontal según $050$ (y $230$);
$\sigma_3$ es horizontal según $140$; $\sigma_2$ es vertical, según la recta en la que se cortan las dos fallas verticales. Cada
falla forma \ang{30} con $\sigma_1$, así que $\ang{45} - \phi/2 =
\ang{30}$ y $\phi = \ang{30}$. (b) Las dos fallas dividen el suelo en cuatro
cuñas; las cuñas que contienen $\sigma_1$, hacia $050$ y hacia $230$, se acercan entre sí, y las que contienen $\sigma_3$ se separan,
hacia $140$ y $320$. Tómese la falla de \omterm{def:b1:geological-maps:attitude}{rumbo} $020$, entre la cuña hacia $050$ (a su
este) y la cuña hacia $320$ (a su oeste): respecto de la cuña occidental, la oriental se
mueve hacia $230$ y se aleja de $320$, y ambas tienen una componente hacia $200$ según la falla. Un observador en la cuña
occidental, mirando al este, ve que la oriental se mueve hacia $200$, a la derecha: la falla de \omterm{def:b1:geological-maps:attitude}{rumbo} $020$ es
\omterm{def:b1:faults-folds-fabrics:fault-types}{dextral}. El mismo argumento, o la simetría respecto de $\sigma_1$, hace que la falla de \omterm{def:b1:geological-maps:attitude}{rumbo} $080$ sea \omterm{def:b1:faults-folds-fabrics:fault-types}{sinistral}.
\end{solution}

\begin{solution}{exo:b1:faults-folds-fabrics:7}
Los vectores normales hacia arriba son
\begin{align*}
  \mathbf n_1 &= (\sin\ang{40}\sin\ang{120}, \sin\ang{40}\cos\ang{120}, \cos\ang{40}) = (0.557, -0.321, 0.766),\\
  \mathbf n_2 &= (\sin\ang{30}\sin\ang{260}, \sin\ang{30}\cos\ang{260}, \cos\ang{30}) = (-0.492, -0.087, 0.866).
\end{align*}
Su producto vectorial es
$(-0.211, -0.859, -0.207)$, que apunta hacia abajo: la \omterm{def:b1:faults-folds-fabrics:fold-elements}{línea de charnela} se inmerge
$\arcsin 0.207/0.909 = \ang{13.1}$ hacia $\ang{180} + \arctan(0.211/0.859) =
\ang{193.8}$. $\mathbf n_1\cdot\mathbf n_2 = 0.417$, el ángulo entre las
normales es \ang{65.3}, y $\iota = \ang{114.7}$: un \omterm{def:b1:faults-folds-fabrics:fold-classes}{pliegue abierto}.
\end{solution}

\begin{solution}{exo:b1:faults-folds-fabrics:8}
$\mathbf n_1 = (0.321, 0.383, 0.866)$ para la estratificación y
$\mathbf n_2 = (0.337, -0.925, 0.174)$ para el \omterm{def:b1:faults-folds-fabrics:foliation}{clivaje}. Su producto vectorial, tomado hacia abajo, da una recta que se inmerge \ang{25.4}
hacia \ang{074.8}. Según \cref{prop:b1:faults-folds-fabrics:intersection}, esta es la dirección de la \omterm{def:b1:faults-folds-fabrics:fold-elements}{línea de charnela}: el \omterm{def:b1:geological-maps:fold}{pliegue} se inmerge unos
\ang{25} hacia el este-nordeste, lo que debe confirmar cualquier otro afloramiento de uno de los \omterm{def:b1:faults-folds-fabrics:fold-elements}{flancos}.
\end{solution}

\begin{solution}{exo:b1:faults-folds-fabrics:9}
(a) Un observador al sur de la falla, mirando al norte, ve la parte norte del dique desplazada hacia el este, a la derecha: la falla es
\omterm{def:b1:faults-folds-fabrics:fault-types}{dextral}. (b) El dique es un plano vertical perpendicular a la falla, así que su desplazamiento según el \omterm{def:b1:geological-maps:attitude}{rumbo} es la componente de
desgarre: \qty{350}{m}, que es todo el deslizamiento si las estrías son horizontales. (c) El deslizamiento forma entonces \ang{20} con
el \omterm{def:b1:geological-maps:attitude}{rumbo} en el plano vertical de la falla: su componente horizontal sigue siendo \qty{350}{m}, así que el \omterm{def:b1:faults-folds-fabrics:fault-elements}{deslizamiento neto} es
$350/\cos\ang{20} = \qty{372.5}{m}$ y su componente vertical $350\tan\ang{20} = \qty{127.4}{m}$. El bloque norte se movió hacia el este, y
las estrías bajan hacia el este, así que también bajó, en \qty{127.4}{m}. Un dique vertical no puede mostrar ese movimiento vertical,
pues deslizarlo a lo largo de sí mismo deja su traza sin cambios.
\end{solution}

\begin{solution}{exo:b1:faults-folds-fabrics:10}
(a) La normal a la falla forma $\ang{90} - \theta$ con $\sigma_1$, así que, según
\cref{prop:b1:stress-strain-rheology:mohr},
$\sigma_n' = \sigma_1'\sin^2\theta + \sigma_3'\cos^2\theta$ y
$\tau = (\sigma_1' - \sigma_3')\sin\theta\cos\theta$. El deslizamiento empieza
cuando $\tau = \mu\sigma_n'$; dividiendo por $\sigma_3'\cos^2\theta$ y escribiendo
$R = \sigma_1'/\sigma_3'$, $(R - 1)\tan\theta = \mu(R\tan^2\theta + 1)$, así que
$R\tan\theta(1 - \mu\tan\theta) = \tan\theta + \mu$, que es la fórmula.
(b) Con $\sigma_1$ vertical, $\theta$ es \ang{90} menos el \omterm{def:b1:geological-maps:attitude}{buzamiento}:
$\theta = \ang{30}$ da $R = 3.12$; $\theta = \ang{45}$, $R = 4.00$;
$\theta = \ang{55}$, $R = 9.92$. (c) El denominador se anula cuando
$\tan\theta = 1/\mu$, $\theta = \ang{59.0}$; más allá, $R$ sería negativo:
ningún cociente de \omterm{def:b1:stress-strain-rheology:lithostatic}{esfuerzos efectivos} positivos puede hacer deslizar la falla, por grande que sea
$\sigma_1'$. Las \omterm{def:b1:faults-folds-fabrics:fault-types}{fallas normales} que buzan menos de \ang{31.0} no pueden deslizar
con $\sigma_1$ vertical y $\mu = 0.6$. (d) El plano de mejor orientación está a
$\ang{45} - \phi/2$ de $\sigma_1$. Para $\mu = 0.6$, $\phi = \ang{31.0}$:
las \omterm{def:b1:faults-folds-fabrics:fault-types}{fallas normales} buzan \ang{60.5}, los \omterm{def:b1:faults-folds-fabrics:fault-types}{cabalgamientos} \ang{29.5}. Para $\mu = 0.85$,
$\phi = \ang{40.4}$: \ang{65.2} y \ang{24.8}.
\end{solution}

\begin{solution}{exo:b1:faults-folds-fabrics:11}
Mirando al norte, el este está a la derecha, y los planos que buzan al oeste suben hacia la derecha. (a) En B la estratificación
buza en el mismo sentido que el \omterm{def:b1:faults-folds-fabrics:foliation}{clivaje}, y más empinada: según \cref{prop:b1:faults-folds-fabrics:cleavage-bedding}(a), B está en un
flanco volcado; A y C, donde la estratificación es menos empinada, están en posición normal. (b) El \omterm{def:b1:faults-folds-fabrics:foliation}{clivaje} sube con \ang{60}; en A y C la
estratificación sube con \ang{40}, y la rotación aguda de uno a otro es en el sentido de las agujas del reloj, a través de
\ang{20}, desplazando la parte superior del \omterm{def:b1:faults-folds-fabrics:foliation}{clivaje} hacia el este; en B la estratificación sube con \ang{80} y la rotación es en sentido
antihorario, a través de \ang{20}, desplazándola hacia el oeste. (c) El sentido cambia entre A y B y de nuevo entre B y C: hay una
\omterm{def:b1:faults-folds-fabrics:fold-elements}{charnela} en cada intervalo. En A, que está en posición normal, la parte superior del \omterm{def:b1:faults-folds-fabrics:foliation}{clivaje} se desplaza hacia el este, hacia el
\omterm{def:b1:geological-maps:fold}{anticlinal} de la parte superior del flanco: la \omterm{def:b1:faults-folds-fabrics:fold-elements}{charnela} entre A y B es un \omterm{def:b1:geological-maps:fold}{anticlinal}, y la de entre B y C un \omterm{def:b1:geological-maps:fold}{sinclinal}.
(d) En la notación de la demostración, los \omterm{def:b1:faults-folds-fabrics:fold-elements}{flancos} en posición normal buzan
$\gamma - \iota/2 = \ang{40}$ y los \omterm{def:b1:faults-folds-fabrics:fold-elements}{flancos} volcados $\gamma + \iota/2 =
\ang{80}$, así que $\gamma = \ang{60}$, como dice el \omterm{def:b1:faults-folds-fabrics:foliation}{clivaje}, e $\iota =
\ang{40}$: \omterm{def:b1:faults-folds-fabrics:fold-classes}{pliegues cerrados}. Los \omterm{def:b1:faults-folds-fabrics:fold-elements}{planos axiales} tienen \omterm{def:b1:geological-maps:attitude}{rumbo} norte y buzan \ang{60} hacia el
oeste, \omterm{def:b1:geological-maps:attitude}{actitud} $270/60$. Los \omterm{def:b1:faults-folds-fabrics:fold-elements}{flancos} cortos han girado más allá de la vertical con sus partes superiores
hacia el este: los \omterm{def:b1:geological-maps:fold}{pliegues} vergen hacia el este, como en la figura del capítulo.
\end{solution}

\begin{solution}{exo:b1:faults-folds-fabrics:12}
(a) En un \omterm{def:b1:geological-maps:fold}{pliegue} en chevrón simétrico cada flanco forma el ángulo $\iota/2$ con el
\omterm{def:b1:faults-folds-fabrics:fold-elements}{plano axial}, que es perpendicular a la dirección del acortamiento; un flanco de longitud $l$ se extiende, por tanto, $l\sin(\iota/2)$
a través de los \omterm{def:b1:geological-maps:fold}{pliegues}, y todo el estrato $L_0\sin(\iota/2)$. (b) $\sin\ang{30} = 0.5$: un acortamiento de
\qty{50}{\percent}. (c) $\sin(\iota/2) = 0.7$ da $\iota = \ang{88.9}$.
(d) $\iota = \ang{156}$ y $\sin\ang{78} = 0.978$: un acortamiento de
\qty{2.2}{\percent}, que es también $1 - \cos\ang{12}$.
\end{solution}

\begin{solution}{pb:b1:faults-folds-fabrics:1}
\textbf{1.} $\mathbf n = (\sin\delta\sin\alpha, \sin\delta\cos\alpha,
\cos\delta)$ da
\[
  \mathbf n_1 = (0.20855, 0.03415, 0.97742), \qquad
  \mathbf n_2 = (-0.20855, 0.03415, 0.97742).
\]

\textbf{2.} $\mathbf n_1\cdot\mathbf n_2 = -0.04349 + 0.00117 + 0.95535 =
0.91302$, y el ángulo entre las normales es \ang{24.1}.

\textbf{3.} Por \cref{prop:b1:faults-folds-fabrics:hinge},
$\iota = \ang{180} - \ang{24.1} = \ang{155.9}$.

\textbf{4.} Entre \ang{170} y \ang{90}: un \omterm{def:b1:faults-folds-fabrics:fold-classes}{pliegue abierto}.

\textbf{5.} Los \omterm{def:b1:faults-folds-fabrics:fold-elements}{flancos} buzan hacia $080.7$ y $279.3$, no exactamente hacia el este y el oeste, porque el \omterm{def:b1:geological-maps:fold}{pliegue} se inmerge: sus
rectas de máxima pendiente no están en el perfil, que solo muestra un \omterm{def:b1:cross-sections:apparent-dip}{buzamiento aparente}. En una sección este--oeste, según
\cref{thm:b1:cross-sections:apparent-dip}, $\tan\beta = \tan\ang{12.2}
\cos\ang{9.3} = 0.2134$, así que $\beta = \ang{12.0}$, y $\ang{180} -
2\times\ang{12.0} = \ang{155.9}$ (el perfil está inclinado \ang{2} respecto de esa sección, lo que no cambia nada con esta precisión).

\textbf{6.} Sus \omterm{def:b1:faults-folds-fabrics:fold-elements}{flancos} son paralelos a los del \omterm{def:b1:geological-maps:fold}{anticlinal}, así que su \omterm{def:b1:faults-folds-fabrics:fold-elements}{ángulo entre flancos} es también \ang{155.9}.

\textbf{7.} $\mathbf n_1\times\mathbf n_2 = (0.03415\times0.97742 -
0.97742\times0.03415,\; 0.97742\times(-0.20855) - 0.20855\times0.97742,\;
0.20855\times0.03415 + 0.03415\times0.20855) = (0, -0.40768, 0.01424)$.
Tomado hacia abajo, $(0, 0.40768, -0.01424)$: la \omterm{def:b1:faults-folds-fabrics:fold-elements}{línea de charnela} tiene tendencia $000$ e \omterm{def:b1:field-lab-methods-1:trend-plunge}{inmersión}
$\arctan(0.01424/0.40768) = \ang{2.0}$ hacia el norte.

\textbf{8.} El \omterm{def:b1:faults-folds-fabrics:fold-elements}{plano axial} contiene la \omterm{def:b1:faults-folds-fabrics:fold-elements}{línea de charnela} y biseca los dos \omterm{def:b1:faults-folds-fabrics:fold-elements}{flancos}; su normal es según
$\mathbf n_1 - \mathbf n_2 = (0.41710, 0, 0)$, horizontal y este--oeste. Es vertical y tiene \omterm{def:b1:geological-maps:attitude}{rumbo} norte--sur: $090/90$. El
\omterm{def:b1:geological-maps:fold}{anticlinal} es erguido.

\textbf{9.} Un \omterm{def:b1:faults-folds-fabrics:fold-classes}{pliegue abierto} y erguido, que se inmerge \ang{2.0} hacia el norte.

\textbf{10.} Según \cref{prop:b1:geological-maps:folds}, los afloramientos de un \omterm{def:b1:geological-maps:fold}{anticlinal} se cierran en la dirección de la
\omterm{def:b1:field-lab-methods-1:trend-plunge}{inmersión}: hacia el norte. Su núcleo contiene las rocas más antiguas: la arenisca de la base de la cobertera y, allí donde el río
ha cortado más hondo, el basamento de esquisto y granito que yace debajo.

\textbf{11.} Axial-plano, el \omterm{def:b1:faults-folds-fabrics:foliation}{clivaje} es vertical y tiene \omterm{def:b1:geological-maps:attitude}{rumbo} norte--sur, $090/90$. Según
\cref{prop:b1:faults-folds-fabrics:intersection}, su \omterm{def:b1:faults-folds-fabrics:lineation}{linación de intersección} con la estratificación es paralela a la \omterm{def:b1:faults-folds-fabrics:fold-elements}{línea de
charnela}: tendencia $000$, \omterm{def:b1:field-lab-methods-1:trend-plunge}{inmersión} \ang{2.0}. Directamente: el plano $080.7/12.2$ corta al plano vertical de \omterm{def:b1:geological-maps:attitude}{rumbo} norte según una
recta de \omterm{def:b1:field-lab-methods-1:trend-plunge}{inmersión} $\arctan(\tan\ang{12.2}
\cos\ang{80.7}) = \ang{2.0}$ hacia $000$.

\textbf{12.} Mirando al norte, el flanco este sube hacia la izquierda con \ang{12.0} y el \omterm{def:b1:faults-folds-fabrics:foliation}{clivaje} es vertical: la rotación aguda del
\omterm{def:b1:faults-folds-fabrics:foliation}{clivaje} a la estratificación, a través de \ang{78.0}, es en sentido antihorario y desplaza la parte superior del \omterm{def:b1:faults-folds-fabrics:foliation}{clivaje} hacia el
oeste. El flanco está en posición normal, así que según \cref{prop:b1:faults-folds-fabrics:cleavage-bedding}(b) la \omterm{def:b1:faults-folds-fabrics:fold-elements}{charnela} del
\omterm{def:b1:geological-maps:fold}{anticlinal} está al oeste, en $x = \qty{3.5}{km}$, como muestra el mapa.

\textbf{13.} El \omterm{def:b1:geological-maps:attitude}{rumbo} es $120 - 90 = 030$; el plano buza hacia $120$, así que el \omterm{def:b1:faults-folds-fabrics:fault-elements}{bloque superior} yace en el lado este-sudeste de la
traza.

\textbf{14.} La falla es normal: $\sigma_1$ era vertical.

\textbf{15.} $\sigma_2$ es horizontal según el \omterm{def:b1:geological-maps:attitude}{rumbo}, $030$--$210$; $\sigma_3$, horizontal y transversal a él, $120$--$300$.

\textbf{16.} $\ang{45} + \phi/2 = \ang{60}$ da $\phi = \ang{30}$, y $\mu = \tan\ang{30} = 0.577$.

\textbf{17.} El conjunto conjugado tiene también \omterm{def:b1:geological-maps:attitude}{rumbo} $030$ y buza \ang{60} en sentido contrario: $300/60$. Los dos conjuntos se
cortan con $\ang{180} - 2\times\ang{60} = \ang{60}$, con la $\sigma_1$ vertical bisecando el ángulo agudo.

\textbf{18.} Con $\mu = 0.85$, $\phi = \ang{40.4}$ y las \omterm{def:b1:faults-folds-fabrics:fault-types}{fallas normales} de mejor orientación buzan \ang{65.2}; la falla de Varden, con
\ang{60}, está a \ang{5.2} de ella. Según \cref{exo:b1:faults-folds-fabrics:10}, desliza con el cociente
$\sigma_1'/\sigma_3' = (1 + 0.85\cot\ang{30})/(1 - 0.85\tan\ang{30}) = 4.85$, frente a $4.68$ del mejor plano: está casi, pero no
exactamente, en la mejor orientación, y volvería a deslizar con un cociente de esfuerzos solo \qty{4}{\percent} mayor.

\textbf{19.} El deslizamiento es hacia abajo según el \omterm{def:b1:geological-maps:attitude}{buzamiento}: $r = \ang{-90}$; el \omterm{def:b1:faults-folds-fabrics:fault-elements}{vector de deslizamiento} tiene tendencia
$120$ e \omterm{def:b1:field-lab-methods-1:trend-plunge}{inmersión} \ang{60}.

\textbf{20.} $(\cos\ang{60}\sin\ang{120}, \cos\ang{60}\cos\ang{120},
-\sin\ang{60}) = (0.4330, -0.2500, -0.8660)$.

\textbf{21.} $e = 200/\tan\ang{60} = \qty{115.5}{m}$.

\textbf{22.} Este: $115.47\sin\ang{120} = \qty{100.0}{m}$; norte:
$115.47\cos\ang{120} = \qty{-57.7}{m}$, es decir, \qty{57.7}{m} hacia el sur.

\textbf{23.} $L = 200/\sin\ang{60} = \qty{230.9}{m}$, y
$200^2 + 115.47^2 = 40\,000 + 13\,333 = 53\,333 = 230.94^2$. El salto y el \omterm{def:b1:cross-sections:throw}{salto horizontal} son las componentes vertical y
horizontal del deslizamiento, que yace en el plano vertical del \omterm{def:b1:geological-maps:attitude}{buzamiento}: son los catetos de un triángulo rectángulo cuya
hipotenusa es el deslizamiento.

\textbf{24.} Redondeando al metro, el \omterm{def:b1:faults-folds-fabrics:fault-elements}{deslizamiento neto} de la falla de Varden es
\textbf{\qty[text-series-to-math=true]{231}{m}}.
\end{solution}
